\documentclass[12pt,a4paper]{article}
\usepackage[utf8]{inputenc}
\usepackage[T1]{fontenc}
\usepackage{polski}
\usepackage[polish]{babel}
\usepackage{geometry}
\geometry{margin=2cm}
\usepackage{graphicx}
\usepackage{tabularx}
\usepackage{booktabs} 
\usepackage{hyperref}
\usepackage{bookmark}
\usepackage{enumitem}
\usepackage{titlesec}
\usepackage{microtype}
\usepackage{varwidth}
\usepackage{multicol}


\title{
    \begin{center}
    \includegraphics[width=\textwidth]{fig/logo.png}
    \vspace{0.25cm}
    \end{center}
    \textbf{\huge Bazy danych} \\
    \vspace{0.5cm}
    \textit{Projekt: Client Relationship Management Application }
    }
\author{Piotr Frąc}
\date{Październik 2026}


\begin{document}
\maketitle
\clearpage


\begin{multicols*}{2}
\subsection*{Wprowadzenie}
CRM - \textbf{Client Relationship Management} jest to oprogramowanie do zarządzania relacjami z klientami, 
które centralizuje dane o kontrahentach i automatyzuje procesy sprzedaży. Gromadzi pełną historie interakcji i ustaleń w jednym miejscu.
 Przejmuje powtarzalne zadania sprzedaży, marketingu i obsługi klienta. Ułatwia monitorowanie celów sprzedażowych oraz zachowań klientów.

\subsection*{Zakres projektu}
\textbf{Musi mieć}
\begin{itemize}
    \item 
\end{itemize}

\textbf{Powinien mieć}
\begin{itemize}
    \item 
\end{itemize}

\textbf{Mogłoby być}
\begin{itemize}
    \item 
\end{itemize}

\textbf{Nie będzie miał}
\begin{itemize}
    \item 
\end{itemize}

\end{multicols*}


\end{document}